% Lab 7 — Path-Based Analysis（原书 PDF p.83–90）

\bichapter{Lab 7：基于路径的分析（PBA）}{Lab 7: Path-Based Analysis}



\begin{objectiveBox}

完成本实验后，你应能够：

\begin{itemize}

  \item 使用 \textbf{path mode PBA} 与 \textbf{exhaustive PBA} 提高 PrimeTime 时序分析精度；

  \item 对违例路径执行交互式路径重算，并比较 GBA 与 PBA 结果；

  \item 生成并解读 PBA 汇总报告。

\end{itemize}

\end{objectiveBox}



\begin{durationBox}

约 \textbf{30} 分钟。

\end{durationBox}



\section{概述}

\label{lab7:overview}



\noindent\textit{Introduction}



本实验将执行 exhaustive PBA（Exhaustive Path-Based Analysis），并在交互模式下对最坏违例路径进行 path mode PBA 重算。



\subsection{相关文件与目录}

\label{lab7:files}



\noindent\textit{Relevant Files and Directories}



本实验全部文件位于主目录下的 \cmd{lab7_pba} 目录中。



\begin{tabularx}{\textwidth}{@{}l X@{}}

\toprule

\textbf{路径 / 文件} & \textbf{说明} \\

\midrule

\cmd{lab7_pba/} & 当前工作目录 \\

\cmd{.synopsys_pt.setup} & 自动读取的 PT setup 文件 \\

\cmd{orca_savesession/} & 已保存会话目录 \\

\cmd{RUN.tcl} & 用于重建 \cmd{orca_savesession} 的运行脚本 \\

\bottomrule

\end{tabularx}



\subsection{答案 / 解答}

\label{lab7:answers-note}



\noindent\textit{Answers / Solutions}



本实验手册在每个 Lab 后半部分给出全部问题的答案与解答。

若答题需要帮助，或希望核对结果，请查阅本 Lab 后部的解答章节。



\section{操作说明}

\label{lab7:instructions}



\noindent\textit{Instructions}



\subsection{任务 1：报告当前设计中的违例路径}

\label{lab7:task1}



\noindent\textit{Task 1. Report violating paths in the current design}



本任务将报告在 \textbf{worst\_slew propagation} 期间计算得到的最坏违例 slack。



\begin{enumerate}

  \item 进入本实验工作目录：

\begin{lstlisting}

unix% cd lab7_pba

\end{lstlisting}



  \item 启动 PrimeTime，并恢复名为 \cmd{orca_savesession} 的已保存会话：

\begin{lstlisting}

unix% pt_shell

pt_shell> restore_session orca_savesession/

\end{lstlisting}



\begin{noteBox}

\begin{itemize}

  \item 如需重建该会话，可使用：\\

        \cmd{pt_shell -f RUN.tcl | tee -i run.log}

  \item 执行 \cmd{RUN.tcl} 时出现的任何 \textbf{PARA-124} 错误，在本课程实验中可安全忽略。

\end{itemize}

\end{noteBox}



  \item 确定 setup 违例数量，并找出到达某端点的最坏违例路径：

\begin{lstlisting}

pt_shell> report_analysis_coverage -check setup \

-status violated

\end{lstlisting}



\begin{questionBox}

\textbf{问题 1.}\ 有多少个 setup 检查处于 violated 状态？\\

\textbf{问题 2.}\ 最坏违例 slack 的大小是多少？对应的端点（endpoint）是什么？

\end{questionBox}

\end{enumerate}



\subsection{任务 2：执行 path mode PBA}

\label{lab7:task2}



\noindent\textit{Task 2. Perform path mode PBA}



本任务将交互式地重算已报告的最坏违例路径的 slack。



\begin{enumerate}

  \item 使用 PBA path 模式确定最坏 setup slack 违例：

\begin{lstlisting}

pt_shell> report_timing -to \

I_ORCA_TOP/I_BLENDER/s4_op2_reg[31]/D -pba_mode path

\end{lstlisting}



  或：

\begin{lstlisting}

pt_shell> set path [get_timing_paths -to \

I_ORCA_TOP/I_BLENDER/s4_op2_reg[31]/D -pba_mode path]

pt_shell> report_timing $path

\end{lstlisting}



\begin{questionBox}

\textbf{问题 3.}\ \cmd{report_timing} 如何表明该路径已被重算？\\

\textbf{问题 4.}\ 交互式 path-based analysis（\opt{-pba\_mode path}）之后，最坏 slack 违例是多少？\\

\textbf{问题 5.}\ 与 GBA slack 相比，PBA slack 是改善、变差还是不变？结果是否符合预期？

\end{questionBox}



\begin{noteBox}

PBA 不会修改 PrimeTime 数据库，因此不能再用 \cmd{report_analysis_coverage} 查看 PBA 结果。

本实验中可使用带 \opt{-pba\_mode} 选项的 \cmd{report_timing} 或 \cmd{get_timing_paths}。

\end{noteBox}

\end{enumerate}



\subsection{任务 3：执行 exhaustive PBA}

\label{lab7:task3}



\noindent\textit{Task 3. Performing Exhaustive PBA}



本任务将在 exhaustive PBA 模式下运行设计分析。



\begin{enumerate}

  \item 找出时钟组 \cmd{SYS_CLK} 中前 10 个 setup 违例及其端点：

\begin{lstlisting}

pt_shell> redirect -tee GBA.rpt { report_timing \

-path_summary -group SYS_CLK -max_paths 10 -nosplit }

\end{lstlisting}



  \item 对设计中的 setup 违例执行 exhaustive PBA，并生成汇总报告：

\begin{lstlisting}

pt_shell> redirect -tee PBA.rpt { report_timing \

-path_summary -group SYS_CLK \

-max_paths 10 -nosplit -pba_mode exhaustive }

\end{lstlisting}



  \item 比较 \cmd{GBA.rpt} 与 \cmd{PBA.rpt}：

\begin{questionBox}

\textbf{问题 6.}\ GBA 与 PBA 各返回多少条违例路径？

\end{questionBox}



\begin{lstlisting}

pt_shell> sizeof_collection [get_timing_paths \

-slack_lesser_than 0 -group SYS_CLK -max_paths 10]

pt_shell> sizeof_collection [get_timing_paths \

-slack_lesser_than 0 -group SYS_CLK -max_paths 10 \

-pba_mode exhaustive]

\end{lstlisting}



\begin{questionBox}

\textbf{问题 7.}\ exhaustive PBA 之后最坏违例的值是多少？对应的时序路径端点名是什么？\\

\textbf{问题 8.}\ 是否出现 \textbf{UITE-480} 警告？

\end{questionBox}

\end{enumerate}



\subsection{任务 4：［可选］生成 GBA 与 PBA 对比汇总报告}

\label{lab7:task4}



\noindent\textit{Task 4. [OPTIONAL Task] Generate GBA vs.\ PBA Summary Reports}



本可选任务将分别用 GBA、PBA path 模式与 PBA exhaustive 模式生成全局、QoR 与约束报告以便对比。



\begin{lstlisting}

redirect -tee GBAsummary.rpt {

  report_global_timing

  report_qor

  report_constraint -all_violators

}

redirect -tee PBApath.rpt {

  report_global_timing -pba_mode path

  report_qor -pba_mode path

  report_constraint -all_violators -pba_mode path

}

redirect -tee PBAexhaust.rpt {

  report_global_timing -pba_mode exhaustive

  report_qor -pba_mode exhaustive

  report_constraint -all_violators -pba_mode exhaustive

}

\end{lstlisting}



\noindent 至此 Lab~7 完成。Day-2 课程结束。



\section{答案 / 解答}

\label{lab7:solutions}



\noindent\textit{Answers / Solutions}



\begin{answerBox}

\textbf{问题 1.}\ 共有 \textbf{23} 个 setup 检查 violated。

\end{answerBox}



\begin{answerBox}

\textbf{问题 2.}\ 最坏违例 slack 为 \textbf{-0.7737}，端点为\\

\cmd{I_ORCA_TOP/I_BLENDER/s4_op2_reg[31]/D}。

\end{answerBox}



\begin{answerBox}

\textbf{问题 3.}\ \textbf{Path Type} 行后出现 \texttt{(recalculated)}。例如：

\begin{lstlisting}

pt_shell> report_timing -to I_ORCA_TOP/I_BLENDER/s4_op2_reg[31]/D \

-pba_mode path

...

Path Group: SYS_CLK

Path Type:  max (recalculated)

\end{lstlisting}

\end{answerBox}



\begin{answerBox}

\textbf{问题 4.}\ 交互式 path-based analysis 之后最坏 slack 为 \textbf{-0.7711}。

\end{answerBox}



\begin{answerBox}

\textbf{问题 5.}\ slack \textbf{有所改善}，符合预期。\\

（路径重算后 slack 只能保持不变或改善。）

\end{answerBox}



\begin{answerBox}

\textbf{问题 6.}\ GBA：\textbf{9} 条违例路径；PBA exhaustive：\textbf{8} 条违例路径。

\end{answerBox}



\begin{answerBox}

\textbf{问题 7.}\ exhaustive PBA 之前最坏违例为 \textbf{-0.7737}，之后为 \textbf{-0.7711}；\\

端点均为 \cmd{I_ORCA_TOP/I_BLENDER/s4_op2_reg[31]/D}。

\end{answerBox}



\begin{answerBox}

\textbf{问题 8.}\ 没有 \textbf{UITE-480} 警告——exhaustive PBA 重算已完成。

\end{answerBox}


